\section{Esperimento: Valutazione delle Prestazioni al Variare dell'Ampiezza della Finestra Temporale}
\label{sec:esperimento_variazione_mu}

\subsection{Obiettivo dell'Esperimento}
L'esperimento implementato nel file \texttt{experiments3.cpp} ha lo scopo di analizzare empiricamente l'efficienza temporale della struttura dati basata su albero temporale (\textbf{Temporal Range Forest / Range Tree}) rispetto a una procedura di ricerca sequenziale non indicizzata (\textbf{Baseline Naive}).
In particolare, l'esperimento valuta il tempo medio di risposta a una query di vicinato temporale con sketch MinHash, osservando come variano le prestazioni al crescere dell'ampiezza della finestra temporale $\mu$.

Dato un nodo target $v \in V$ e un intervallo temporale $I = [t_{\text{start}}, t_{\text{end}}]$ di ampiezza $\mu = t_{\text{end}} - t_{\text{start}} + 1$, l'obiettivo di entrambe le metodologie è calcolare la signature MinHash $S(v, I) = \operatorname{Sketch}(N^I(v))$ associata al vicinato temporale:
\[
N^I(v) = \{ u \in V \mid \exists t' \in I : (u, v, t') \in E \}.
\]

\subsection{Funzionamento e Metodologia}
L'esperimento si articola in quattro fasi sequenziali:

\begin{enumerate}
    \item \textbf{Generazione del Grafo Temporale Sintetico:}
    Viene generato uno scenario di test ad alta densità temporale mediante la funzione \texttt{generate\_synthetic\_graph}. I parametri utilizzati sono:
    \begin{itemize}
        \item Numero totale di vertici: $|V| = 70\,000$;
        \item Grado medio: $\bar{d} = 350$, per un volume complessivo di archi temporali pari a $|E| = |V| \cdot \bar{d} = 24\,500\,000$;
        \item Orizzonte temporale discreto: $\mathcal{T} = \{1, \ldots, T_{\max}\}$, con $T_{\max} = 5481$.
    \end{itemize}
    Gli archi $(u, v, t)$ vengono generati selezionando casualmente coppie di nodi distinti $u \neq v$ e assegnando un timestamp uniforme $t \in [1, T_{\max}]$. Viene infine estratto casualmente un nodo query target $v \in V$.

    \item \textbf{Costruzione della Struttura Pre-indicizzata (TRF):}
    Viene istanziato \emph{una sola volta all'inizio} il \texttt{RangeTree} associato al nodo target $v$:
    \begin{itemize}
        \item Per ogni istante attivo $t \in \Lambda(v)$, viene calcolato lo sketch foglia MinHash con dimensione fissata a $k = 128$ permutazioni.
        \item I nodi interni memorizzano l'unione dei vicinati dei sottoalberi tramite l'operazione di \texttt{merge} tra sketch MinHash (prendendo il minimo componente per componente, con costo $O(k)$).
        \item Il tempo di costruzione viene ammortizzato una tantum prima dell'avvio della fase di interrogazione.
    \end{itemize}

    \item \textbf{Campionamento Dinamico delle Query per Differenti Ampiezze $\mu$:}
    L'esperimento itera su una sequenza di ampiezze temporali crescenti:
    \[
    \mu \in \{50, 100, 200, 500, 1000, 2000, 5000\}.
    \]
    Per ciascun valore di $\mu \le T_{\max}$:
    \begin{itemize}
        \item Si calcola il numero teorico massimo di intervalli ammissibili:
        \[
        \text{max\_possible\_intervals} = T_{\max} - \mu + 1.
        \]
        \item Si generano $Q = 100$ query campionando uniformemente $t_{\text{start}} \in [1, T_{\max} - \mu + 1]$ e fissando $t_{\text{end}} = t_{\text{start}} + \mu - 1$.
    \end{itemize}

    \item \textbf{Confronto Algoritmico a Due Vie:}
    Per ogni batch di $Q$ query, vengono misurati i tempi di esecuzione ad alta risoluzione (\texttt{std::chrono::high\_resolution\_clock}):
    \begin{itemize}
        \item \textbf{Temporal Range Forest (\texttt{tree->query}):} Risolve la query identificando i nodi canonici dell'albero che partizionano l'intervallo $[t_{\text{start}}, t_{\text{end}}]$ ed eseguendo il \texttt{merge} dei rispettivi sketch. Il costo asintotico atteso è:
        \[
        O(k \log |\Lambda(v)|).
        \]
        \item \textbf{Approccio Naive (\texttt{naive\_query}):} Esegue una scansione lineare dell'intero insieme degli archi $E$, filtra quelli incidenti su $v$ con timestamp $t \in [t_{\text{start}}, t_{\text{end}}]$ e ricalcola da zero la signature MinHash (costo $O(|E| + k \cdot |N^I(v)|)$).
    \end{itemize}
\end{enumerate}

\subsection{Descrizione dell'Output}
L'esperimento produce due flussi di output:

\subsubsection{1. Output su File CSV (\texttt{results\_interval2.csv})}
I dati numerici aggregati vengono esportati in formato tabellare CSV per permettere la generazione di grafici e l'analisi statistica. La struttura delle colonne è la seguente:
\begin{itemize}
    \item \textbf{\texttt{mu}}: ampiezza della finestra temporale $\mu$ considerata.
    \item \textbf{\texttt{max\_possible\_intervals}}: numero massimo di finestre distinte di durata $\mu$ collocabili nell'orizzonte temporale $[1, T_{\max}]$, calcolato come $T_{\max} - \mu + 1$.
    \item \textbf{\texttt{trf\_time\_ns}}: tempo medio di esecuzione per query impiegato dalla Temporal Range Forest, espresso in nanosecondi (calcolato come media su $Q = 100$ istanze).
    \item \textbf{\texttt{naive\_time\_ns}}: tempo medio per query impiegato dall'algoritmo Naive per ricostruire lo sketch da zero, espresso in nanosecondi.
\end{itemize}

\subsubsection{2. Log a Terminale (Console)}
Durante l'esecuzione, il programma stampa a video:
\begin{itemize}
    \item La conferma della generazione del grafo e l'identificativo del nodo target selezionato casualmente;
    \item Lo stato di completamento della costruzione una tantum del \texttt{RangeTree};
    \item Una riga riassuntiva per ciascun valore di $\mu$ che riporta il confronto diretto tra il tempo della TRF e il tempo della baseline naive.
\end{itemize}

\subsection{Interpretazione Teorica dei Risultati}
L'analisi empirica evidenzia due comportamenti strutturali fondamentali:
\begin{enumerate}
    \item \textbf{Scalabilità della TRF rispetto all'ampiezza dell'intervallo:}
    Mentre la baseline naive subisce un incremento del tempo di elaborazione all'aumentare di $\mu$ (poiché una finestra più ampia intercetta un numero maggiore di eventi temporali da scansionare ed inserire nello sketch), la Temporal Range Forest mantiene tempi di risposta logaritmici $O(k \log |\Lambda(v)|)$ stabili, poiché l'intervallo viene decomposto in al più $O(\log |\Lambda(v)|)$ nodi canonici indipendentemente dalla durata assoluta di $\mu$.
    
    \item \textbf{Fattore di Accelerazione (Speedup):}
    Il rapporto
    \[
    \text{Speedup} = \frac{\text{naive\_time\_ns}}{\text{trf\_time\_ns}}
    \]
    quantifica il beneficio dell'indicizzazione gerarchica, dimostrando ordini di grandezza di vantaggio rispetto alla scansione esaustiva della timeline, specialmente per finestre temporali ampie dove il ricalcolo da zero risulterebbe proibitivo.
\end{enumerate}